\section{Goals and Research Questions}
\label{intro:goals}

In recent years, large language models (LLMs) have demonstrated unprecedented capabilities in natural language processing tasks, including translation, question answering, reasoning-based problem solving, and automated code generation.
Their success has fueled a wide range of popular applications—including chatbots, search engines, and code generation agents—that rely on LLMs as a fundamental component of their backend systems. As a result, LLM inference has become one of the most requested software-as-a-service (SaaS) workloads.

Numerous solutions have been proposed to address the challenge of accelerating LLM inference at scale. Currently, graphics processing units (GPUs) represent the most widely adopted approach, while domain-specific architectures (DSAs)—particularly dataflow AI accelerators—are emerging as a promising alternative \cite{chen2020survey, koilia2024hardware, silvano2025survey, kachris2025survey}.
The primary distinction between GPUs and dataflow accelerators depends on their approach to parallel computation \cite{kundu2025comparison}:
\begin{itemize}
    \item \textbf{GPUs} employ the single-instruction, multiple-thread (SIMT) model to execute thousands of threads in parallel, thereby accelerating the matrix and vector operations characteristic of deep learning workloads.
    \item \textbf{Dataflow accelerators} stream data through a network of compute units in a highly pipelined, parallel, and often asynchronously scheduled manner. This design significantly reduces data movement.
\end{itemize}

Beyond pursuing higher performance, accelerator designs increasingly prioritise energy efficiency to ensure sustainable large-scale deployment.
According to the International Energy Agency (IEA), the energy consumption of accelerated servers is projected to increase by 30\% annually from 2024 to 2030, primarily due to AI-related demand \cite{iea2025ai}.
Notably, companies such as NVIDIA and AWS report that 80–90\% of their AI-related energy consumption arises from model serving (inference) rather than training \cite{patterson2021carbon}.
Consequently, minimising the energy consumption of LLM inference is crucial for ensuring the economic and environmental sustainability of these technologies.

Recent studies have examined LLM inference across different hardware accelerators. Most prior work focuses on NVIDIA GPUs \cite{fernandez2025energy, wilhelm2025beyond, wilkins2024offline, stojkovic2025dynamollm, maliakel2025investigating, kakolyris2024slo, stojkovic2024towards}, while others also include AMD and Intel GPUs \cite{reddi2020mlperf, LLMInfbench} and dataflow accelerators \cite{emani2024toward, Emani2022, LLMInfbench}.
However, key gaps remain: there is limited cross-vendor evaluation directly comparing modern GPUs, with most studies emphasising performance while overlooking energy efficiency; and there is a lack of investigation into dataflow accelerators.

Bridging these gaps is challenging because LLM inference performance depends on many interrelated factors beyond the choice of hardware accelerators. We identify four key challenges:

\begin{enumerate}
    \item \textbf{Model characteristics}: The parameter count and architectural design of an LLM, such as the use of Mixture-of-Experts layers, strongly affect both computational demands and memory footprint.

    \item \textbf{Batch size}: This parameter introduces complex trade-offs, as it directly influences the balance between computation and memory access patterns, often increasing memory pressure and impacting overall performance.

    \item \textbf{Quantisation}: Modern GPUs frequently employ quantisation to reduce memory requirements and improve computational efficiency; however, its effects vary significantly across architectures, highlighting the need for explicit, architecture-aware evaluations.

    \item \textbf{Multi-GPU serving}: Modern LLMs often require multiple accelerators for inference. A variety of parallelisation strategies—such as tensor, pipeline, or data parallelism—can be employed, each exhibiting distinct, hardware-dependent effects on performance scaling.
\end{enumerate}

Together, these factors underscore the intrinsic complexity of systematically analysing performance and energy efficiency in LLM inference across diverse accelerator platforms.


With this thesis, we aim to provide a comprehensive evaluation of LLM inference performance and energy consumption across multiple AI accelerators, including datacenter-level GPUs from Nvidia, AMD, and Intel, as well as dataflow AI accelerators from SambaNova and Cerebras.
We evaluate fourteen open-source LLMs with different architectures and scales, and analyse the effects of model size, batch size, quantisation, and deployment configurations ranging from single- and multi-GPU setups to dataflow-accelerator deployments. This comprehensive analysis reveals trade-offs among these factors, performance, and energy efficiency, offering deeper insight into how models and hardware interact under realistic inference workloads.